\section{Chaos in the Noosphere}

\paragraph{Three bodies}


\textbf{Shankara}, following the teaching of the Vedanta, describes the three bodies that constitute man.

\begin{enumerate}


\item Gross body: physical senses

\item Subtle body, or the mind/psyche: emotions, feelings, attractions, aversions, etc

\item Causal body, or the intellect: thoughts, images
\end{enumerate}


Together, they make up the phenomenal world, i.e., the world as experienced, or \emph{Jagat}. Note, particularly, that what we commonly regard as our ``inner life'' is actually part of the phenomenal world.

The intellect is called the causal body because it is the cause of our experiences in the subtle and gross bodies. Few people are aware of that process occurring in their own consciousnesses. That is because they regard the outer world as real, independent, and objective, but their thoughts as subjective, and actually the result of events in the physical world, not the other way around.

Shankara describes the activity of the causal body as a waking dream, analogous to the sleeping dream state. Liberation comes from knowledge of the Self. This Self, or Atma, transcends the causal body, and is unaffected by the vagaries and divagations of the causal, subtle, and gross bodies. Atman is consciousness itself, and that is what brings the activities of the three bodies into awareness.

Unfortunately, we do not understand our identity, certainly not as Atman, but we identify as a false finite self, a creation of the causal body. Due to ignorance, the causal body creates an illusory world that we believe we inhabit. When and if we wake up to our true identity, that world dissipates, just as a nighttime dream does when we awaken from sleep. It is pointless to look for that ``ignorance'', since it is not real, as it is only a privation; one cannot find ignorance, any more than one can find the darkness while walking around with a torch (as they say).

\paragraph{Chaos}


It is difficult to convince anyone that he is ignorant; since the identification with the ego is so strong, it seems undeniably certain to him. Nevertheless, by pointing out the ill effects of ignorance, perhaps some will be prompted to pursue things a little more deeply.

First of all, it is clear that the experiences of the gross body, i.e., everything experienced through the senses, is remarkably stable (apart from special cases of malfunctioning sensory organs). In other words, we all go outside and agree that the sky is blue and the grass is green. If I ask three people in the know for directions to the nearest convenience store, they will pretty much agree and I will arrive at my destination. The relative reliability of the gross world will be of value in the beginning exercises for the achieving Gnosis.

This stability cannot be explained if the causal body is randomly creating sensory experiences. Shankara elaborates a system of five elements (space, fire, air, water, earth), five pranas, ten organs, manas, and buddhi. These are subtle states that determine the gross material world. So the awakening from the limited ego to knowledge of the true self, involves a change in the perspective of the observer, not in the outer world. \textbf{Swami Chinmayananda} explains:

\begin{quotex}
If and when after that awakening there is an outward cognition of any experience, such a Mahatman of Self-realization cannot but see the same equipments of experiences, which were before his own, singing the Eternal Song of Life.
\end{quotex}


The situation is reversed for the subtle body, which appears to be entirely subjective. Clearly, there are emotional states that are reflected in bodily gestures. But for the most part, our inner states are invisible to each other; there are anxieties, fears, concerns, desires, etc., that dominate our minds. That is because of the identification with the ego, which results in a constant state of negativity. This makes interpersonal relationships so difficult: each party experiences the world filtered through these inner states and is unable to take into account the inner state of the other. To do so, is an art; first one must become detached from his own identifications.

\paragraph{Noosphere}


The activity of the causal body in those identified with the ego is also private and subjective, yet has an interpersonal element due to the faculty of speech. Speech then communicates systems of thought or ideologies. We can call the totality of public thought the ``Noosphere'', in accordance with common usage. However, the Noosphere is not necessarily a boon, as it is actually in chaos.

With no way to arbitrate incompatible ideologies, they multiply and ramify. Since the identification with one's own ideologies is so strong, any threat to it is misconstrued as a threat to one's very self. Imagine, in our other examples, that every time you asserted that the ``sky is blue'', you would find yourself embroiled in arguments with your neighbors, and sometimes even very bitter arguments. Then if you asked three friends for directions, two of them would give you totally incompatible directions and the third would try to convince you not to go to the convenience store at all. Life would be difficult and tedious; you would eventually carve out enough partial knowledge to function. Furthermore, you would form alliances with those holding similar ideologies. That is the chaos in the Noosphere.

The obvious question that arises is whether there is a way to settle ideological disputes, analogous to the way the gross world is organized. Some will turn to scientific positivism, which is the study of the orderliness of the gross world. That provides some sense of consensus, since mathematicians and scientists can agree on their thought systems; or at least, they have a recognized and mutually accepted way to resolve such differences, relying on experiments in and observations of the gross world. But this comes at the heavy cost of limiting what counts as knowledge; the rest of life, e.g., philosophy, political opinions, etc., does not lend itself to scientific methods.

Keep in mind that when we say we ``know'' something in the physical world, that really means that we intuit its form in the intellectual soul, even if we are unaware of doing so. So is there an analogous intuition of the thought world? Yes, there is, in fact, and that is the revelation of Tradition. \textbf{Valentin Tomberg} calls that ``the book'', so we have the progression of thought to speech to the book. Exoterically, this revelation seems to be just another ideology competing for space in the Noosphere. The result is that the original revelation gets distorted and ramifies into multiple versions. A Traditional society is organized around its originary revelatory impulse, which maintains the populace in the awareness of that revelation. Rites, prayers, sacraments, scriptures, and so on, will awaken a little intuition in people so that they can maintain that faith. But this system is fragile, depending on the fuller understanding by the spiritual leaders.

The intuition in relation to the Intellect is not of the forms of objects, but rather of the Self. This this is ultimately One, there can be only one revelation. The method described by Shankara is remarkably similar to what Tomberg calls the birth of the Logos from the Spirit and the Virgin. Shankara uses the analogy of water reflecting the light of the moon. If the Intellect is agitated by attachments, desires, pleasures, pains, likes, dislikes, and so on, the moon's reflection is fragmented. However, if the water is still, the moon's reflection in the water is a unity.

Shankara concludes with a threefold plan to rid the ego of such impurities:

\begin{enumerate}


\item Listen to the truth of Scriptures, i.e., hear the Word

\item Reason from those truths

\item Deep contemplation on what has been heard and reasoned out
\end{enumerate}


Shankara does not proceed beyond the realization of the self; at that point, the task is complete, the realizer is immortal, and there is no longer any concern with the manifest world. This is unlike the Buddhist who claim that the Bodhisattva will return to the world. It also differs from the Western Tradition, which sees the end in the fulfillment of the Person, not in his merging into the Absolute.

Reference: \emph{Atma Bodha}, by \textbf{Shankara}, commentary by \textbf{Swami Chinmayananda}


Next: whether Tomberg and Solovyov offer anything further


\flrightit{Posted on 2014-03-24 by Cologero}

\begin{center}* * *\end{center}

\begin{footnotesize}
\begin{sffamily}
\texttt{Michel on 2014-03-25 at 06:56 said:}

``Don't be afraid of the man behind the curtain''.

Seeing the vast trickery in the world's sophists and in myself gave life to that statement. Except it has been pretty dangerous along the way sometimes. I wonder if anyone here has also experienced any danger in the psychic world, and how they dealt with it. That man seems very powerful before the curtain is torn and he is shown for what he is.

\hfill

\texttt{scardanelli on 2014-03-25 at 10:17 said:}

``A Traditional society is organized around its originary revelatory impulse, which maintains the populace in the awareness of that revelation. Rites, prayers, sacraments, scriptures, and so on, will awaken a little intuition in people so that they can maintain that faith. But this system is fragile, depending on the fuller understanding by the spiritual leaders.''

This, then , underlies the importance of the ``elite'' of a society- the leaven that infuses the whole. The elite is only elite insofar as it has the capacity for revelation, for the intuition of the true self rather than merely intellectual understanding. It seems that this ``post'' is perhaps that of the hermit in Tomberg... that of the prophet in distinction to the posts of emperor and pope.

``... such a Mahatman of Self-realization cannot but see the same matter equipments of experiences... ''

Cologero, could you clear up this passage? I don't quite understand it. Is there a typo maybe?

\hfill

\texttt{Tom on 2014-03-25 at 11:41 said:}

Dear Cologero.

Always read your posts with great interest.

Just one question regarding the '' three bodies'' : I thought that thoughts and images were an aspect of the subtle body which corresponds to the dream state? The causal body corresponds to the deep sleep state .

\hfill

\texttt{Frater M on 2014-03-25 at 12:25 said:}

Reminds one of Paul to the Ephesians ``For our struggle is not against flesh and blood, but against the rulers, against the powers, against the world forces of this darkness, against the spiritual forces of wickedness in the heavenly places''.

\hfill

\texttt{andros on 2014-03-25 at 16:26 said:}

On dealing with danger: many years ago, while abroad, a malevolent presence entered my mind, then my room. Turning over in my bed and facing the wall only made the presence stronger, it's incursion more frightening. Not knowing what to do, I brought to mind the people I love: fear vanished and I recognised sickness. I was rising to help the presence when my neighbour started knocking frantically on the wall.

Love brought me through.

\hfill

\texttt{Cologero on 2014-03-25 at 19:36 said:}

\begin{quotationx}\footnotesize

Chinmayandanda: We are not limited individuality as we take ourselves to be. There is a limitless region of experience beyond what is known at present,and this is shut out from us because of our self-centred existence. When the realisation of the true nature of the Self comes to one--even as mere theoretical knowledge--he, who was till then in a perpetual state of fear and sorrow, seems to wake up from his own limited existence and rediscovers himself to be the Blissful Self which is All-pervading and Eternal.

If and when thereafter there is an outward cognition of any experience, such a Mahatman of Self-realisation cannot but see the same matter as equipments of experiences, which were before his own, singing the Eternal Song of Life.
\end{quotationx}


By equipments he means the physical, mental, and intellectual . It seems to me that he is saying that after self-realisation the world is seen as it is, and not through the illusions we create.

\hfill

\texttt{Cologero on 2014-03-25 at 20:04 said:}

Tom, the Swami refers to ``the world of objects cognised for enjoyment as forms, smells, etc., through the sense-organs; as feelings and emotions experienced by the mind; and as ideas and ideologies lived by the intellect.''

I'm glad you pointed this out, because after posting it, I wanted to go back and insert a step prior to ``thought, speech, book''; I was a little quick equating the intellect to the causal body because of the way the commentary was written. Yes, the thought arises in the mind, but is caused by the ignorance in the causal body (hence its name). It does not change the general point that ideologies arise out of ignorance and they dictate the way we live.

\hfill

\texttt{Michel on 2014-03-26 at 12:13 said:}

@andros

Thanks for sharing that. I hope that in continuing with your self-realization you have been able to progress through the subtle world overcoming the vast trickery therein. Of course the danger is in the trickery itself, and in the indefinite forms it is capable of assuming, be it the disgusting and terrifying or the lovely and beautiful. A tasty meal should look as disgusting as a bowl of dung; and a bowl of dung should look as delightful as a tasty meal. All is Aoristos after all.

\end{sffamily}
\end{footnotesize}

